\chapter{2012年大纲卷（文）}

\section{单选题}

\begin{problem}
已知集合 \(A = \{x \mid x\text{ 是平行四边形}\}\)，\(B = \{x \mid x\text{ 是矩形}\}\)，\(C = \{x \mid x\text{ 是正方形}\}\)，\(D = \{x \mid x\text{ 是菱形}\}\)，则（\quad）

\choices
    {\(A \subseteq B\)}
    {\(C \subseteq B\)}
    {\(D \subseteq C\)}
    {\(A \subseteq D\)}
\end{problem}

\begin{answer}
B
\end{answer}

\begin{solution}
正方形的两组对边分别平行且四个角都是直角，所以正方形一定是矩形，即 \(C\subseteq B\). 一般的平行四边形不一定有直角，菱形也不一定是正方形，也存在不是菱形的平行四边形，因此其余包含关系不成立，故选 B.
\end{solution}

\begin{problem}
函数 \( y = \sqrt{x + 1} \) （\(x \geqslant -1\)）的反函数为（\quad）

\choices
    {\( y = x^{2} - 1 \)（\(x \geqslant 0\)）}
    {\( y = x^{2} - 1 \)（\(x \geqslant 1\)）}
    {\( y = x^{2} + 1 \)（\(x \geqslant 0\)）}
    {\( y = x^{2} + 1 \)（\(x \geqslant 1\)）}
\end{problem}

\begin{answer}
A
\end{answer}

\begin{solution}
由 \(y=\sqrt{x+1}\) 可知 \(y\geqslant0\). 交换 \(x\)，\(y\) 得 \(x=\sqrt{y+1}\)，两边平方可得 \(y=x^{2}-1\)，并且由原函数的值域得到反函数的定义域 \(x\geqslant0\). 因此反函数为选项 A，故选 A.
\end{solution}

\begin{problem}
若函数 \( f(x) = \sin \frac{x + \varphi}{3} \)（\(\varphi \in [0, 2\pi]\)）是偶函数，则 \( \varphi = \)（\quad）

\choices
    {\(\frac{\pi}{2}\)}
    {\(\frac{2\pi}{3}\)}
    {\(\frac{3\pi}{2}\)}
    {\(\frac{5\pi}{3}\)}
\end{problem}

\begin{answer}
C
\end{answer}

\begin{solution}
令 \(u=\frac{x}{3}\)，则
\[
f(x)=\sin\left(u+\frac{\varphi}{3}\right)
\]
\[
f(-x)=\sin\left(-u+\frac{\varphi}{3}\right)
\]
函数为偶函数要求两式对任意 \(u\) 相等. 利用正弦的和差公式，
\[
f(x)-f(-x)=2\cos\frac{\varphi}{3}\sin u
\]
要使其恒为零，必须有 \(\cos\frac{\varphi}{3}=0\). 又 \(\varphi\in[0,2\pi]\)，所以 \(\frac{\varphi}{3}=\frac{\pi}{2}\)，从而 \(\varphi=\frac{3\pi}{2}\)，故选 C.
此时 \(f(x)=\cos\frac{x}{3}\)，确为偶函数，所以该值满足题设.
\end{solution}

\begin{problem}
已知 \(\alpha\) 为第二象限角，\(\sin \alpha = \frac{3}{5}\)，则 \(\sin 2\alpha =\)（\quad）

\choices
    {\(-\frac{24}{25}\)}
    {\(-\frac{12}{25}\)}
    {\(\frac{12}{25}\)}
    {\(\frac{24}{25}\)}
\end{problem}

\begin{answer}
A
\end{answer}

\begin{solution}
因为 \(\alpha\) 是第二象限角，所以
\[
\cos\alpha=-\sqrt{1-\sin^{2}\alpha}=-\sqrt{1-\frac{9}{25}}=-\frac45
\]
于是
\[
\sin2\alpha=2\sin\alpha\cos\alpha
=2\times\frac35\times\left(-\frac45\right)=-\frac{24}{25}
\]
故选 A.
\end{solution}

\begin{problem}
椭圆的中心在原点，焦距为 \(4\)，一条准线为 \(x = -4\)，则该椭圆的方程为（\quad）

\choices
    {\(\frac{x^2}{16} +\frac{y^2}{12} = 1\)}
    {\(\frac{x^2}{12} +\frac{y^2}{8} = 1\)}
    {\(\frac{x^2}{8} +\frac{y^2}{4} = 1\)}
    {\(\frac{x^2}{12} +\frac{y^2}{4} = 1\)}
\end{problem}

\begin{answer}
C
\end{answer}

\begin{solution}
设椭圆的半长轴、半短轴和半焦距分别为 \(a\)，\(b\)，\(c\)，离心率为 \(e=\frac{c}{a}\). 焦距为 \(4\) 给出 \(2c=4\)，所以 \(c=2\). 由于准线为 \(x=-4\)，椭圆的长轴在 \(x\) 轴上，准线到原点的距离满足
\[
\frac{a}{e}=\frac{a^{2}}{c}=4
\]
因此 \(a^{2}=8\)，再由 \(b^{2}=a^{2}-c^{2}\) 得 \(b^{2}=4\). 所以椭圆方程为
\[
\frac{x^{2}}8+\frac{y^{2}}4=1
\]
故选 C.
\end{solution}

\begin{problem}
已知数列 \(\{a_{n}\}\) 的前 \(n\) 项和为 \(S_{n}\)，\(a_{1} = 1\)，\(S_{n} = 2a_{n+1}\)，则 \(S_{n} =\)（\quad）

\choices
    {\(2^{n - 1}\)}
    {\(\left(\frac{3}{2}\right)^{n - 1}\)}
    {\(\left(\frac{2}{3}\right)^{n - 1}\)}
    {\(\frac{1}{2^{n - 1}}\)}
\end{problem}

\begin{answer}
B
\end{answer}

\begin{solution}
当 \(n\geqslant2\) 时，\(a_n=S_n-S_{n-1}\)，而由已知关系有 \(a_n=\frac12S_{n-1}\). 所以
\[
S_n=S_{n-1}+a_n=\frac32S_{n-1}
\]
又 \(S_1=a_1=1\)，递推得到
\[
S_n=\left(\frac32\right)^{n-1}
\]
故选 B.
\end{solution}

\begin{problem}
\(6\) 位选手依次演讲，其中选手甲不在第一个也不在最后一个演讲，则不同的演讲次序共有（\quad）

\choices
    {\(240\) 种}
    {\(360\) 种}
    {\(480\) 种}
    {\(720\) 种}
\end{problem}

\begin{answer}
C
\end{answer}

\begin{solution}
选手甲可以安排在第 \(2\)，\(3\)，\(4\)，\(5\) 个位置，共有 \(4\) 种选择. 确定甲的位置后，其余 \(5\) 位选手可以任意排列，有 \(5!=120\) 种. 因此不同演讲次序共有
\[
4\times5!=4\times120=480
\]
种，故选 C.
\end{solution}

\begin{problem}
已知正四棱柱 \(ABCD - A_{1}B_{1}C_{1}D_{1}\) 中，\(AB = 2\)，\(CC_{1} = 2\sqrt{2}\)，\(E\) 为 \(CC_{1}\) 的中点，则直线 \(AC_{1}\) 到平面 \(BED\) 的距离为（\quad）

\choices
    {\(2\)}
    {\(\sqrt{3}\)}
    {\(\sqrt{2}\)}
    {\(1\)}
\end{problem}

\begin{answer}
D
\end{answer}

\begin{solution}
取正四棱柱的坐标为 \(A(0,0,0)\)，\(B(2,0,0)\)，\(C(2,2,0)\)，\(D(0,2,0)\)，
\(C_1(2,2,2\sqrt2)\)，\(E(2,2,\sqrt2)\).
平面 \(BED\) 的两个方向向量可取 \(\overrightarrow{DB}=(2,-2,0)\) 和 \(\overrightarrow{DE}=(2,0,\sqrt2)\)，其法向量可取 \((\sqrt2,\sqrt2,-2)\). 故平面方程为
\[
\sqrt2x+\sqrt2y-2z-2\sqrt2=0
\]
直线 \(AC_1\) 的方向向量为 \((2,2,2\sqrt2)\)，与该法向量的内积为零，所以直线 \(AC_1\) 平行于平面 \(BED\). 因此所求距离等于点 \(A\) 到平面的距离：
\[
d=\frac{|-2\sqrt2|}{\sqrt{(\sqrt2)^2+(\sqrt2)^2+(-2)^2}}=\frac{2\sqrt2}{2\sqrt2}=1
\]
故选 D.
\end{solution}

\begin{problem}
\(\triangle ABC\) 中，\(AB\) 边的高为 \(CD\)，若 \(\overrightarrow{CB} = \bs{a}\)，\(\overrightarrow{CA} = \bs{b}\)，\(\bs{a} \cdot \bs{b} = 0\)，\(|\bs{a}| = 1\)，\(|\bs{b}| = 2\)，则 \(\overrightarrow{AD} =\)（\quad）

\choices
    {\(\frac{1}{3}\bs{a} - \frac{1}{3}\bs{b}\)}
    {\(\frac{2}{3}\bs{a} - \frac{2}{3}\bs{b}\)}
    {\(\frac{3}{5}\bs{a} - \frac{3}{5}\bs{b}\)}
    {\(\frac{4}{5}\bs{a} - \frac{4}{5}\bs{b}\)}
\end{problem}

\begin{answer}
D
\end{answer}

\begin{solution}
以 \(C\) 为原点，则 \(B\) 的位置向量为 \(\bs a\)，\(A\) 的位置向量为 \(\bs b\). 直线 \(AB\) 的方向向量为 \(\bs a-\bs b\). 设垂足 \(D\) 满足
\[
\overrightarrow{CD}=\bs b+t(\bs a-\bs b)
\]
因为 \(CD\perp AB\)，有
\[
\bigl[\bs b+t(\bs a-\bs b)\bigr]\cdot(\bs a-\bs b)=0
\]
利用 \(\bs a\cdot\bs b=0\)，\(|\bs a|=1\)，\(|\bs b|=2\)，可得
\[
-4+5t=0
\]
所以 \(t=\frac45\). 因此
\[
\overrightarrow{AD}=t\overrightarrow{AB}=\frac45(\bs a-\bs b)=\frac45\bs a-\frac45\bs b
\]
故选 D.
\end{solution}

\begin{problem}
已知 \(F_{1}\)，\(F_{2}\) 为双曲线 \(C: x^{2} - y^{2} = 2\) 的左、右焦点，点 \(P\) 在 \(C\) 上，\(|PF_{1}| = 2|PF_{2}|\)，则 \(\cos \angle F_{1}PF_{2} =\)（\quad）

\choices
    {\(\frac{1}{4}\)}
    {\(\frac{3}{5}\)}
    {\(\frac{3}{4}\)}
    {\(\frac{4}{5}\)}
\end{problem}

\begin{answer}
C
\end{answer}

\begin{solution}
双曲线可写为 \(\frac{x^2}{2}-\frac{y^2}{2}=1\)，所以 \(a=\sqrt2\)，焦半距 \(c=2\)，从而 \(|F_1F_2|=4\). 由于 \(|PF_1|>|PF_2|\)，点 \(P\) 在右支上，双曲线定义给出
\[
|PF_1|-|PF_2|=2a=2\sqrt2
\]
设 \(|PF_2|=t\)，由 \(|PF_1|=2|PF_2|\) 得 \(t=2\sqrt2\)，即 \(|PF_1|=4\sqrt2\). 在三角形 \(F_1PF_2\) 中使用余弦定理，
\[
\cos\angle F_1PF_2
=\frac{(4\sqrt2)^2+(2\sqrt2)^2-4^2}{2\cdot4\sqrt2\cdot2\sqrt2}
=\frac{24}{32}=\frac34
\]
故选 C.
\end{solution}

\begin{problem}
已知 \(x = \ln \pi\)，\(y = \log_5 2\)，\(z = \e^{-\frac{1}{2}}\)，则（\quad）

\choices
    {\(x < y < z\)}
    {\(z < x < y\)}
    {\(z < y < x\)}
    {\(y < z < x\)}
\end{problem}

\begin{answer}
D
\end{answer}

\begin{solution}
因为 \(5>1\)，所以对数函数 \(\log_5x\) 单调递增. 又 \(2<\sqrt5\)，故
\[
y=\log_52<\frac12
\]
由 \(e<4\) 得 \(z=e^{-1/2}>\frac12\)，所以 \(y<z\). 此外 \(e<3<\pi\)，从而 \(x=\ln\pi>1\)，而 \(z<1\)，故 \(z<x\). 综上 \(y<z<x\)，故选 D.
\end{solution}

\begin{problem}
正方形 \(ABCD\) 的边长为 \(1\)，点 \(E\) 在边 \(AB\) 上，点 \(F\) 在边 \(BC\) 上，\(AE = BF = \frac{1}{3}\). 动点 \(P\) 从 \(E\) 出发沿直线向 \(F\) 运动，每当碰到正方形的边时反弹，反弹时反射角等于入射角，当点 \(P\) 第一次碰到 \(E\) 时，\(P\) 与正方形的边碰撞的次数为（\quad）

\choices
    {\(8\)}
    {\(6\)}
    {\(4\)}
    {\(3\)}
\end{problem}

\begin{answer}
B
\end{answer}

\begin{solution}
取正方形为 \(0\leqslant x\leqslant1\)，\(0\leqslant y\leqslant1\)，令 \(A=(0,0)\). 则 \(E=(\frac13,0)\)，\(F=(1,\frac13)\)，运动方向的斜率为 \(\frac12\). 将正方形关于边反射展开后，运动轨迹仍是一条直线，可取其方向向量为 \((2,1)\).

从 \(E\) 出发，轨迹第一次回到 \(E\) 的镜像点时，纵坐标增加 \(2\)，横坐标增加 \(4\)，即到达展开图中的 \(E'=(\frac{13}{3},2)\). 在这段展开轨迹中依次穿过竖直网格线 \(x=1\)，\(2\)，\(3\)，\(4\)，穿过水平网格线 \(y=1\)，最后到达 \(E'\). 前五次是四次竖直边和一次水平边，最后回到原正方形的 \(E\) 时又碰到一次边，共碰撞 \(6\) 次，故选 B.
\end{solution}

\section{填空题}

\begin{problem}
\(\left(x + \frac{1}{2x}\right)^8\) 的展开式中 \(x^2\) 的系数为\fillinblank{}.
\end{problem}

\begin{answer}
\(7\)
\end{answer}

\begin{solution}
展开式的通项为
\[
\mathrm C_8^k x^{8-k}\left(\frac{1}{2x}\right)^k
=\mathrm C_8^k\frac{1}{2^k}x^{8-2k}
\]
要得到 \(x^2\)，必须满足 \(8-2k=2\)，所以 \(k=3\). 因此所求系数为
\[
\mathrm C_8^3\frac{1}{2^3}=\frac{56}{8}=7
\]
\end{solution}

\begin{problem}
若 \(x\)，\(y\) 满足约束条件 \(\begin{cases}
x - y + 1 \geqslant 0,\\
x + y - 3 \leqslant 0,\\
x + 3y - 3 \geqslant 0,
\end{cases}\) 则 \(z = 3x - y\) 的最小值为\fillinblank{}.
\end{problem}

\begin{answer}
\(-1\)
\end{answer}

\begin{solution}
三个约束分别化为
\[
\begin{cases}
y\leqslant x+1,\\
y\leqslant3-x,\\
y\geqslant1-\frac{x}{3}.
\end{cases}
\]
可行域是由这三条直线围成的三角形. 求交点得
\(y=x+1\)，\(\ y=3-x\Longrightarrow (x,y)=(1,2)\)
\(y=x+1\)，\(\ y=1-\frac{x}{3}\Longrightarrow (x,y)=(0,1)\)
\(y=3-x\)，\(\ y=1-\frac{x}{3}\Longrightarrow (x,y)=(3,0)\).
目标函数 \(z=3x-y\) 为线性函数，最小值在顶点取得. 代入三个顶点分别得到
 \(z(1,2)=1\)，\(z(0,1)=-1\)，\(z(3,0)=9\).
所以最小值为 \(-1\).
\end{solution}

\begin{problem}
当函数 \( y = \sin x - \sqrt{3}\cos x \)（\(0 \leqslant x < 2\pi\)）取得最大值时，\( x = \)\fillinblank{}.
\end{problem}

\begin{answer}
\(\frac{5\pi}{6}\)
\end{answer}

\begin{solution}
利用辅助角公式，
\[
\sin x-\sqrt3\cos x
=2\left(\frac12\sin x-\frac{\sqrt3}{2}\cos x\right)
=2\sin\left(x-\frac{\pi}{3}\right)
\]
函数最大值为 \(2\)，且取最大值时
\[
x-\frac{\pi}{3}=\frac{\pi}{2}+2k\pi
\]
结合 \(0\leqslant x<2\pi\)，唯一可取的 \(x\) 为 \(\frac{5\pi}{6}\).
\end{solution}

\begin{problem}
已知正方体 \(ABCD - A_{1}B_{1}C_{1}D_{1}\) 中，\(E\)，\(F\) 分别为 \(BB_{1}\)，\(CC_{1}\) 的中点，那么异面直线 \(AE\) 与 \(D_{1}F\) 所成角的余弦值为\fillinblank{}.
\end{problem}

\begin{answer}
\(\frac{3}{5}\)
\end{answer}

\begin{solution}
设正方体棱长为 \(1\)，取坐标 \(A(0,0,0)\)，\(B(1,0,0)\)，\(C(1,1,0)\)，
\(D_1(0,1,1)\). 则 \(E\left(1,0,\frac12\right)\)，
\(F\left(1,1,\frac12\right)\).
可取两条异面直线的方向向量
\[
\overrightarrow{AE}=\left(1,0,\frac12\right)
\]
\[
\overrightarrow{D_1F}=\left(1,0,-\frac12\right)
\]
它们的内积为 \(\frac34\)，两向量长度均为 \(\frac{\sqrt5}{2}\)，所以所成锐角的余弦为
\[
\frac{\frac34}{\frac{\sqrt5}{2}\cdot\frac{\sqrt5}{2}}=\frac35
\]
\end{solution}

\section{解答题}

\begin{problem}
\(\triangle ABC\) 的内角 \(A\)，\(B\)，\(C\) 成等差数列，其对边 \(a\)，\(b\)，\(c\) 满足 \(2b^{2} = 3ac\)，求 \(A\).
\end{problem}

\begin{solution}
由 \(A\)，\(B\)，\(C\) 成等差数列，得 \(A+C=2B\). 又 \(A+B+C=\pi\)，所以 \(3B=\pi\)，即
\[
B=\frac{\pi}{3}
\]
由余弦定理，
\[
b^2=a^2+c^2-2ac\cos B=a^2+c^2-ac
\]
代入 \(2b^2=3ac\)，得
\[
2a^2+2c^2-2ac=3ac
\]
即
\[
2a^2-5ac+2c^2=0
\]
因为 \(a\)，\(c>0\)，令 \(t=\frac{a}{c}\)，则
\[
2t^2-5t+2=0
\]
所以 \(t=2\) 或 \(t=\frac12\).

当 \(a=2c\) 时，\(b^2=3c^2\)，由余弦定理
\[
\cos A=\frac{b^2+c^2-a^2}{2bc}
=\frac{3c^2+c^2-4c^2}{2\sqrt3c^2}=0
\]
故 \(A=\frac{\pi}{2}\).

当 \(a=\frac12c\) 时，\(b^2=\frac34c^2\)，于是
\[
\cos A=\frac{\frac34c^2+c^2-\frac14c^2}
{2\cdot\frac{\sqrt3}{2}c^2}
=\frac{\sqrt3}{2}
\]
由于 \(0<A<\pi\)，故 \(A=\frac{\pi}{6}\).
\end{solution}

\begin{problem}
已知数列 \( \{a_{n}\} \) 中， \( a_{1}=1 \)，前 \(n\) 项和 \( S_{n}=\frac{n+2}{3}a_{n} \).
\begin{enumerate}
\item 求 \(a_{2}\)，\(a_{3}\).
\item 求 \( \{a_{n}\} \) 的通项公式.
\end{enumerate}
\end{problem}

\begin{solution}
\begin{enumerate}
\item 由 \(S_2=1+a_2=\frac43a_2\)，得 \(a_2=3\). 再由 \(S_3=1+3+a_3=\frac53a_3\)，得 \(a_3=6\).
\item 当 \(n\geqslant2\) 时，利用 \(S_n-S_{n-1}=a_n\)，有
\[
\frac{n+2}{3}a_n-\frac{n+1}{3}a_{n-1}=a_n
\]
整理得
\[
(n-1)a_n=(n+1)a_{n-1}
\]
所以
\[
a_n=\frac{n+1}{n-1}a_{n-1}
\]
因此
\[
a_n=a_1\prod_{k=2}^{n}\frac{k+1}{k-1}
=\frac{3\cdot4\cdots(n+1)}{1\cdot2\cdots(n-1)}
=\frac{n(n+1)}{2}
\]
\end{enumerate}
\end{solution}

\begin{problem}
如图，四棱锥 \(P - ABCD\) 中，底面 \(ABCD\) 为菱形，\(PA \perp\) 底面 \(ABCD\)，\(AC = 2\sqrt{2}\)，\(PA = 2\)，\(E\) 是 \(PC\) 上的一点，\(PE = 2EC\).
\begin{enumerate}
\item 证明：\(PC \perp\) 平面 \(BED\).
\item 设二面角 \(A - PB - C\) 为 \(90^{\circ}\)，求 \(PD\) 与平面 \(PBC\) 所成角的大小.
\end{enumerate}
\bitmapfigure[max width=0.42\linewidth]{img/2012/syllabus_liberal/q19_fig1.png}
\end{problem}

\begin{solution}
\begin{enumerate}
\item 以 \(A\) 为原点建立空间直角坐标系. 由于菱形的两条对角线互相垂直且互相平分，可设
\(A(0,0,0)\)，\(C(2\sqrt2,0,0)\)，\(B(\sqrt2,-b,0)\)，\(D(\sqrt2,b,0)\)，
其中 \(b>0\). 由 \(PA\perp\) 底面且 \(PA=2\)，可取 \(P(0,0,2)\). 由 \(PE=2EC\)，点 \(E\) 将 \(PC\) 按 \(2:1\) 分割，所以
\[
E=P+\frac23(C-P)=\left(\frac{4\sqrt2}{3},0,\frac23\right)
\]
于是
\[
\overrightarrow{PC}=(2\sqrt2,0,-2)
\]
\[
\overrightarrow{BE}=\left(\frac{\sqrt2}{3},b,\frac23\right)
\]
\[
\overrightarrow{DE}=\left(\frac{\sqrt2}{3},-b,\frac23\right)
\]
直接计算得
\[
\overrightarrow{PC}\cdot\overrightarrow{BE}
=2\sqrt2\cdot\frac{\sqrt2}{3}-2\cdot\frac23=0
\]
\[
\overrightarrow{PC}\cdot\overrightarrow{DE}
=2\sqrt2\cdot\frac{\sqrt2}{3}-2\cdot\frac23=0
\]
直线 \(BE\) 与 \(DE\) 相交于 \(E\)，且都在平面 \(BED\) 内，所以 \(PC\perp\) 平面 \(BED\).

\item 平面 \(PAB\) 的一个法向量可取
\[
\bs{p}=(b,\sqrt2,0)
\]
平面 \(PBC\) 的一个法向量可取
\[
\bs{n}=\left(1,-\frac{\sqrt2}{b},\sqrt2\right)
\]
因为它分别与 \(\overrightarrow{PC}\) 和 \(\overrightarrow{BE}\) 垂直. 二面角 \(A-PB-C\) 为 \(90^{\circ}\)，即平面 \(PAB\perp\) 平面 \(PBC\)，所以
\[
\bs{p}\cdot\bs{n}
=b-\frac{2}{b}=0
\]
由 \(b>0\) 得 \(b=\sqrt2\)，此时可取平面 \(PBC\) 的法向量为
\[
\bs{n}=(1,-1,\sqrt2)
\]
又
\[
\overrightarrow{DP}=P-D=(-\sqrt2,-\sqrt2,2)
\]
因此
\[
\frac{\left|\overrightarrow{DP}\cdot\bs{n}\right|}
{|\overrightarrow{DP}|\,|\bs{n}|}
=\frac{2\sqrt2}{(2\sqrt2)\cdot2}=\frac12
\]
若 \(PD\) 与平面 \(PBC\) 所成的角为 \(\theta\)，则上式为 \(\sin\theta\)，所以
\[
\sin\theta=\frac12
\]
所以 \(\theta=30^{\circ}\).
\end{enumerate}
\end{solution}

\begin{problem}
乒乓球比赛规则规定：一局比赛，双方比分在\(10\)平前，一方连续发球\(2\)次后，对方再连续发球\(2\)次，依次轮换. 每次发球，胜方得\(1\)分，负方得\(0\)分. 设在甲，乙的比赛中，每次发球，发球方得\(1\)分的概率为\(0.6\)，各次发球的胜负结果相互独立. 甲，乙的一局比赛中，甲先发球.
\begin{enumerate}
\item 求开始第 \(4\) 次发球时，甲，乙的比分为 \(1\) 比 \(2\) 的概率.
\item 求开始第 \(5\) 次发球时，甲得分领先的概率.
\end{enumerate}
\end{problem}

\begin{solution}
\begin{enumerate}
\item 前两次由甲发球，甲每次得分的概率为 \(0.6\)；第三次由乙发球，甲得分的概率为 \(0.4\). 记前两次中甲得分的次数为 \(X\)，则
\[
P(X=1)=\mathrm C_2^1(0.6)(0.4)=0.48
\]
\[
P(X=0)=(0.4)^2=0.16
\]
三次发球后比分为 \(1:2\) 有两种互斥情形：前两次甲得 \(1\) 分且第三次乙得分，或前两次甲没有得分且第三次甲得分. 因此所求概率为
\[
0.48\times0.6+0.16\times0.4=0.288+0.064=0.352
\]
\item 设前两次发球中甲得分次数为 \(X\)，后两次发球中甲得分次数为 \(Y\). 则
\[
P(X=1)=0.48
\]
\[
P(X=2)=0.6^2=0.36
\]
\[
P(Y=1)=\mathrm C_2^1(0.4)(0.6)=0.48
\]
\[
P(Y=2)=0.4^2=0.16
\]
开始第五次发球时已经完成四次发球，甲得分领先等价于 \(X+Y>2\). 所以
\[
P(X+Y>2)
=P(X=1,Y=2)+P(X=2,Y=1)+P(X=2,Y=2)
\]
\[
=0.48\times0.16+0.36\times0.48+0.36\times0.16
=0.0768+0.1728+0.0576=0.3072
\]
\end{enumerate}
\end{solution}

\begin{problem}
已知函数 \( f(x) = \frac{1}{3} x^3 + x^2 + ax \).
\begin{enumerate}
\item 讨论 \( f(x) \) 的单调性.
\item 设 \( f(x) \) 有两个极值点 \(x_{1}\)，\(x_{2}\)，若过两点 \((x_{1}, f(x_{1}))\)，\((x_{2}, f(x_{2}))\) 的直线 \( l \) 与 \( x \) 轴的交点在曲线 \( f(x) \) 上，求 \( a \) 的值.
\end{enumerate}
\end{problem}

\begin{solution}
\begin{enumerate}
\item
\[
f'(x)=x^2+2x+a=(x+1)^2+a-1
\]
当 \(a\geqslant1\) 时，\(f'(x)\geqslant0\) 对一切 \(x\) 成立，且仅在 \(a=1\) 时于 \(x=-1\) 处为零，因此 \(f(x)\) 在 \(\R\) 上单调递增.

当 \(a<1\) 时，令 \(u=-1-\sqrt{1-a}\)，\(v=-1+\sqrt{1-a}\)，
则 \(u<v\) 且 \(f'(x)\) 在 \((-\infty,u)\) 和 \((v,+\infty)\) 上为正，在 \((u,v)\) 上为负. 因此 \(f(x)\) 在 \((-\infty,u)\)，\((v,+\infty)\) 上单调递增，在 \((u,v)\) 上单调递减.

\item 由“有两个极值点”知 \(a<1\). 设 \(u=-1-s\)，\(v=-1+s\)，
\(s=\sqrt{1-a}>0\)，
则 \(u\)，\(v\) 是 \(f'(x)=0\) 的两个根，且 \(u+v=-2\)，\(uv=a\). 在极值点处 \(a=-u^2-2u\)，故
\[
f(u)=-\frac23u^3-u^2
\]
\[
f(v)=-\frac23v^3-v^2
\]
过两极值点的直线斜率为
\[
m=\frac{f(v)-f(u)}{v-u}
=\frac13(u^2+uv+v^2)+u+v+a
=\frac{2(a-1)}3=-\frac{2s^2}{3}
\]
由 \(u=-1-s\) 及 \(f(u)=\frac{(1+s)^2(2s-1)}3\)，该直线与 \(x\) 轴的交点横坐标为
\[
t=u-\frac{f(u)}m
=-1-s+\frac{(1+s)^2(2s-1)}{2s^2}
=\frac12-\frac{1}{2s^2}
=-\frac{a}{2(1-a)}
\]
由于 \(f(x)=x\left(\frac13x^2+x+a\right)\)，条件“交点在曲线 \(y=f(x)\) 上”等价于
\(t=0\) 或 \(t^2+3t+3a=0\).
当 \(t=0\) 时，\(a=0\). 当 \(t\ne0\) 时，代入 \(t=-\frac{a}{2(1-a)}\)，得
\[
12a^2-17a+6=0
\]
所以 \(a=\frac23\) 或 \(a=\frac34\).
这三个值都小于 \(1\)，均满足有两个极值点，故 \(a=0\)，\(a=\frac23\) 或 \(a=\frac34\).
\end{enumerate}
\end{solution}

\section{选考题：解答题}

\begin{problem}
已知抛物线 \(C: y = (x + 1)^2\) 与圆 \(M: (x - 1)^2 + (y - \frac{1}{2})^2 = r^2\)（\(r > 0\)）有一个公共点 \(A\)，且在 \(A\) 处两曲线的切线为同一直线 \(l\).
\begin{enumerate}
\item 求 \(r\).
\item 设 \(m\)，\(n\) 是异于 \(l\) 且与 \(C\) 及 \(M\) 都相切的两条直线，\(m\)，\(n\) 的交点为 \(D\)，求 \(D\) 到 \(l\) 的距离.
\end{enumerate}
\end{problem}

\begin{solution}
\begin{enumerate}
\item 设公共点 \(A=(x,(x+1)^2)\). 抛物线在 \(A\) 处的切向量可取 \((1,2(x+1))\). 圆心为 \(M=(1,\frac12)\)，而圆的半径 \(AM\) 垂直于切线，所以
\[
(x-1)+2(x+1)\left((x+1)^2-\frac12\right)=0
\]
令 \(t=x+1\)，上式化为
\[
t-2+2t\left(t^2-\frac12\right)=2t^3-2=0
\]
因此 \(t=1\)，即 \(A=(0,1)\). 所以
\[
r=AM=\sqrt{(0-1)^2+\left(1-\frac12\right)^2}
=\frac{\sqrt5}{2}
\]
此时抛物线在 \(A\) 处的切线斜率为 \(2\)，故
\[
l:y=2x+1
\]

\item 设抛物线在点 \(Q=(a,(a+1)^2)\) 处的切线为 \(L_a\). 由导数或切线公式，
\[
L_a:y=2(a+1)x+1-a^2
\]
它与圆相切，当且仅当圆心 \(M=(1,\frac12)\) 到该直线的距离等于 \(r=\frac{\sqrt5}{2}\)，即
\[
\frac{\left|2(a+1)-\frac12+1-a^2\right|}
{\sqrt{4(a+1)^2+1}}=\frac{\sqrt5}{2}
\]
两边平方并化简，得
\[
a^2(a^2-4a-6)=0
\]
其中 \(a=0\) 对应已知切线 \(l\)，其余两个切点参数为
\[
a_1=2-\sqrt{10}
\]
\[
a_2=2+\sqrt{10}
\]
相应的两条公切线 \(m\)，\(n\) 为
\[
y=2(a_1+1)x+1-a_1^2
\]
\[
y=2(a_2+1)x+1-a_2^2
\]
由 \(a_i^2-4a_i-6=0\)，将 \(x=2\) 代入任一条直线的右端得到
\[
4(a_i+1)+1-a_i^2=4a_i+5-a_i^2=-1
\]
所以 \(m\)，\(n\) 的交点为 \(D=(2,-1)\). 直线 \(l\) 的一般式为 \(2x-y+1=0\)，故
\[
D\text{ 到 }l\text{ 的距离}
=\frac{|2\cdot2-(-1)+1|}{\sqrt{2^2+(-1)^2}}
=\frac6{\sqrt5}
=\frac{6\sqrt5}{5}
\]
\end{enumerate}
\end{solution}
